% Source: PDF pages 27--32.
\section{Bivariate generating functions}
A generating function tracks objects by size. To track other properties, we use more variables.
\begin{definition}
A \textit{combinatorial class with a parameter} $(C,|\cdot|,\chi)$ consists of
\begin{itemize}
\item a set of objects $C$;
\item a size function $|\cdot|\colon C\to\N$ such that there are finitely many objects of each size;
\item a parameter function $\chi\colon C\to\N$.
\end{itemize}
Its \textit{bivariate generating function} is
\[
C(u,z)=\sum_{\sigma\in C}u^{\chi(\sigma)}z^{|\sigma|}
=\sum_{n\ge0}\left(\sum_{k\ge0}c_{k,n}u^k\right)z^n,
\]
where $c_{k,n}$ is the number of objects in $C$ with size $n$ and parameter value $k$.
\end{definition}
Setting $u=1$ forgets the parameter:
\[
C(1,z)=\sum_{\sigma\in C}z^{|\sigma|}=C(z).
\]
\begin{example}
Let $B$ be the class of binary strings with size given by length and parameter given by the number of $0$s. Then
\begin{align*}
B(u,z)&=\sum_{n\ge0}\left(\sum_{k\ge0}\binom nk u^k\right)z^n\\
&=\sum_{n\ge0}(1+u)^nz^n\\
&=\sum_{n\ge0}\bigl((1+u)z\bigr)^n
=\frac1{1-(1+u)z}.
\end{align*}
This calculation starts with the counting series. We also want a way to pass directly from classes to generating functions.
\end{example}

\subsection{Markers}
To build bivariate generating functions, add a \textit{marker} $\mu$ to the constructions.
\begin{itemize}
\item $\mu$ has size $0$ and parameter value $1$, so its generating function is $z^0u^1=u$.
\item $\mathcal E$ has size $0$ and parameter value $0$, so its generating function is $1$.
\item $\mathcal Z$ has size $1$ and parameter value $0$, so its generating function is $z$.
\end{itemize}
The marker keeps track of the parameter without contributing to size.
\begin{example}
For binary strings enumerated by length and number of zeros,
\[
B=\operatorname{SEQ}(\mathcal Z_0\times\mu+\mathcal Z_1).
\]
Thus
\[
B(u,z)=\frac1{1-uz-z}=\frac1{1-(1+u)z}.
\]
\end{example}
\begin{example}
Let $C$ be the class of integer compositions, tracking the number of parts equal to $1$. Its specification is
\[
C=\operatorname{SEQ}\bigl(\mathcal Z\times\mu
  +\mathcal Z^2\times\operatorname{SEQ}(\mathcal Z)\bigr).
\]
The first term is a marked part of size $1$; the second represents parts of size at least $2$. Hence
\begin{align*}
C(u,z)&=\frac1{1-\left(zu+\frac{z^2}{1-z}\right)}\\
&=\frac{1-z}{1-(1+u)z+(u-1)z^2}.
\end{align*}
This is a two-variable rational generating function.
\end{example}
\begin{example}
Find the generating function for rooted planar binary trees enumerated by their number of vertices and number of leaves.\par\medskip
Let $T$ be the class of nonempty binary trees. Then
\[
T=\mathcal Z+\mathcal Z\times\mathcal E\times T
 +\mathcal Z\times T\times\mathcal E+\mathcal Z\times T\times T.
\]
These cases correspond to no children, an empty left subtree and nonempty right subtree, a nonempty left subtree and empty right subtree, and two nonempty subtrees.
Marking leaves gives
\[
T=\mu\times\mathcal Z+\mathcal Z\times\mathcal E\times T
 +\mathcal Z\times T\times\mathcal E+\mathcal Z\times T\times T,
\]
and therefore
\[
T(u,z)=uz+2zT(u,z)+zT(u,z)^2.
\]
Solving the quadratic gives
\[
T(u,z)=\frac{1-2z-\sqrt{1-4z+4(1-u)z^2}}{2z}.
\]
% TODO: verify the handwritten z^2 coefficient in the radical on PDF page 29.
% It is unclear/inconsistent with the displayed quadratic. The coefficient
% here is obtained directly from (1-2z)^2-4uz^2.
The other sign introduces a $1/z$ term, so the minus sign is the formal power series solution. Setting $u=1$ recovers the ordinary generating function for nonempty binary trees. The first terms are
\[
T(u,z)=uz+2uz^2+(u^2+4u)z^3+\cdots.
\]
\begin{center}
\begin{tikzpicture}[x=.7cm,y=.7cm,
  vertex/.style={circle,fill=noteink,inner sep=1.5pt},
  leaf/.style={rectangle,draw=notelink,fill=white,inner sep=2.3pt}]
% One vertex, two trees with two vertices, five trees with three vertices.
\node[leaf] at (0,0) {};
\node[vertex] (a) at (1.5,0) {}; \node[leaf] (b) at (1.1,-.65) {}; \draw (a)--(b);
\node[vertex] (c) at (2.8,0) {}; \node[leaf] (d) at (3.2,-.65) {}; \draw (c)--(d);
\node[vertex] (a) at (4.6,0) {};
\node[leaf] (b) at (4.2,-.65) {}; \node[leaf] (c) at (5,-.65) {};
\draw (b)--(a)--(c);
\foreach \x/\s/\t in {6.2/1/1,7.7/1/-1,9.3/-1/1,10.8/-1/-1}{
  \node[vertex] (a) at (\x,0) {};
  \node[vertex] (b) at (\x+.4*\s,-.65) {};
  \node[leaf] (c) at (\x+.4*\s+.4*\t,-1.3) {};
  \draw (a)--(b)--(c);
}
\node[below] at (0,-1.65) {$uz$};
\node[below] at (2.15,-1.65) {$2uz^2$};
\node[below] at (7.7,-1.65) {$(u^2+4u)z^3$};
\end{tikzpicture}
\end{center}

\end{example}

\subsection{Expected values}
Given $\displaystyle C(u,z)=\sum_{n\ge0}\left(\sum_{k\ge0}c_{k,n}u^k\right)z^n$, its formal partial derivative with respect to $u$ is
\[
\frac{\partial}{\partial u}C(u,z)=C_u(u,z)
=\sum_{n\ge0}\left(\sum_{k\ge1}k c_{k,n}u^{k-1}\right)z^n.
\]
The \textit{expected value}, or average, of the parameter $\chi\colon C\to\N$ among the objects of size $n$ is
\begin{align*}
\mathbb E_n[\chi]
&=\sum_{k\ge0}k\,\mathbb P\bigl(\chi(\sigma)=k\mid |\sigma|=n\bigr)\\
&=\sum_{k\ge0}k\,
\frac{\#\{\sigma: \chi(\sigma)=k,\ |\sigma|=n\}}
{\#\{\sigma:|\sigma|=n\}}\\
&=\sum_{k\ge0}k\frac{c_{k,n}}{c_n}.
\end{align*}
Here an object of size $n$ is chosen uniformly at random.
\begin{proposition}
For each $n$ with $c_n\ne0$,
\[
\mathbb E_n[\chi]=\frac{[z^n]C_u(1,z)}{[z^n]C(1,z)}.
\]
\end{proposition}
\begin{proof}
Since
\[
[z^n]C_u(u,z)=\sum_{k\ge1}k c_{k,n}u^{k-1},
\qquad [z^n]C_u(1,z)=\sum_{k\ge0}k c_{k,n},
\]
dividing by $[z^n]C(1,z)=c_n$ gives the expected value.
\end{proof}
\begin{example}
Find the average number of $0$s among binary strings of length $n$.\par\medskip
We know that $\displaystyle B(u,z)=\frac1{1-(1+u)z}$, so
\[
B_u(1,z)=\left.\frac{z}{(1-(1+u)z)^2}\right|_{u=1}
=\frac z{(1-2z)^2}.
\]
For $n\ge1$, the negative binomial theorem gives
\begin{align*}
[z^n]\frac z{(1-2z)^2}
&=[z^{n-1}](1-2z)^{-2}\\
&=2^{n-1}[z^{n-1}](1-z)^{-2}\\
&=2^{n-1}\binom n{n-1}=n2^{n-1}.
\end{align*}
Thus the expected value is $\displaystyle\frac{n2^{n-1}}{2^n}=\frac n2$. For $n=0$, it is also $0$.
\end{example}
\begin{example}
The generating function for compositions enumerated by size and number of parts equal to $1$ is
\[
C(u,z)=\frac{1-z}{(1-z)(1-uz)-z^2},
\qquad [z^n]C(1,z)=2^{n-1}\quad(n\ge1).
\]
Furthermore,
\begin{align*}
C_u(1,z)
&=\left.\frac{(1-z)^2z}{\bigl((1-z)(1-uz)-z^2\bigr)^2}\right|_{u=1}\\
&=\frac{(1-z)^2z}{(1-2z)^2}.
\end{align*}
For $n\ge3$,
\begin{align*}
[z^n]C_u(1,z)
&=[z^n]\frac{z-2z^2+z^3}{(1-2z)^2}\\
&=[z^{n-1}](1-2z)^{-2}
 -2[z^{n-2}](1-2z)^{-2}
 +[z^{n-3}](1-2z)^{-2}\\
&=2^{n-2}+n2^{n-3}.
\end{align*}
Thus, for $n\ge3$, the average number of parts equal to $1$ is
\[
\frac{2^{n-2}+n2^{n-3}}{2^{n-1}}=\frac{n+2}{4}.
\]
For $n=0,1,2$, the averages are respectively $0$, $1$, and $1$.
\end{example}

\subsection{Variance and concentration}
The \textit{variance} of the parameter $\chi$ on objects of size $n$ is
\[
\operatorname{Var}_n[\chi]=\mathbb E_n[\chi^2]-\mathbb E_n[\chi]^2,
\qquad
\mathbb E_n[\chi^2]=\sum_{k\ge0}k^2\frac{c_{k,n}}{c_n}.
\]
The second derivative contributes $k(k-1)$, so adding the first derivative gives the factor $k^2$.
% Source annotation says the second derivative gives k^2 directly;
% the displayed variance formula correctly also includes C_u.
\begin{proposition}
For $c_n\ne0$,
\[
\operatorname{Var}_n[\chi]
=\frac{[z^n]\bigl(C_{uu}(1,z)+C_u(1,z)\bigr)}{[z^n]C(1,z)}
-\left(\frac{[z^n]C_u(1,z)}{[z^n]C(1,z)}\right)^2.
\]
\end{proposition}
\begin{proposition}
Suppose $\mathbb E_n[\chi]\ne0$ for all sufficiently large $n$, and
\[
\lim_{n\to\infty}\frac{\operatorname{Var}_n[\chi]}{\mathbb E_n[\chi]^2}=0.
\]
Then, for every $\varepsilon>0$, the probability that a uniformly random object $\sigma\in C$ of size $n$ satisfies
\[
(1-\varepsilon)\mathbb E_n[\chi]\le\chi(\sigma)
\le(1+\varepsilon)\mathbb E_n[\chi]
\]
tends to $1$ as $n\to\infty$.
\end{proposition}
% Source has (1-epsilon) at both ends of this interval; the upper bound
% must be (1+epsilon), as also indicated by the accompanying sketch.
\begin{center}
\begin{tikzpicture}[x=1.1cm,y=1cm,>=Stealth]
\draw[->] (-.4,0)--(5.7,0);
\draw[->] (0,-.2)--(0,2.1);
\draw[thick,notelink] plot[smooth] coordinates {(0,.02) (1,.08) (1.7,.2) (2.1,.6) (2.4,1.4) (2.7,1.85) (3,1.3) (3.4,.45) (4,.2) (5.4,.06)};
\draw[dashed] (2.7,0)--(2.7,2);
\draw[dashed] (2.05,0)--(2.05,.65);
\draw[dashed] (3.35,0)--(3.35,.65);
\node[below,font=\small] at (2.7,-.05) {$\mathbb E_n[\chi]$};
\draw[<->] (2.05,-.65)--(3.35,-.65);
\node[below,font=\small] at (2.7,-.65) {relative interval of width $2\varepsilon$};
\end{tikzpicture}
\end{center}

